\documentclass[../../../../main.tex]{subfiles}
\begin{document}

{}
　$\triangle OAB$の重心$G$を通る直線が，辺$OA$，$OB$とそれぞれ辺上の点$P$，$Q$で交わっているとする．
$\overrightarrow{OP}=h\overrightarrow{OA}$，$\overrightarrow{OQ}=k\overrightarrow{OB}$とし，
$\triangle OAB$，$\triangle OPQ$の面積をそれぞれ$S$，$T$とすれば，次の関係が成り立つことを示せ．
\begin{description}
\item[(i)]$\displaystyle\frac{1}{h}+\frac{1}{k}=3$
\item[(ii)]$\displaystyle\frac{4}{9}S\leqq T \leqq\frac{1}{2}S$
\end{description}

\end{document}
